

\illustration{0.3}{alley}{A foggy alley. \href{https://pixabay.com/de/photos/park-in-der-nacht-dunkle-stra\%C3\%9Fe-4270765/}{Photo by Henryk Niestr\'{o}j}\vspace{-0.5cm}}%
One cold and foggy night, you walk down a shady alley.
There should be a lamp every few metres but none of them seem to work, and in this night, not even the moon enlightens your path.
Alone and in the dark, you wonder:
``Even if there was a working lamp somewhere, how much would it lighten my way?''.
Now, back at home, you want to calculate this.

The alley can be modelled as a line with a length of $n$ metres.
The fog has a uniform density and reduces the light of a lamp by a factor of $1-p$ every metre.
The brightness at one point is the sum of the light that reaches this point from every lamp.
You want to calculate this brightness at some points after placing some lamps.

